% !TEX program = xelatex
% 统计刊物风格论文模板（中文正文，Times 西文，三线表，定理环境）
% 用法：复制为 workspace/paper.tex，替换 <...> 占位内容，然后
%   xelatex paper.tex && xelatex paper.tex
% 或调用工具 ht_compile_latex(tex="workspace/paper.tex")
\documentclass[10pt,a4paper]{article}

% ---- 中文与版面（xeCJK 直接按文件挂字体，绕开 ctex 的 fontset 探测）----
\usepackage{xeCJK}
\setCJKmainfont[Path=C:/Windows/Fonts/,FontIndex=0]{simsun.ttc}
% 注意：xeCJK 只挂 main font；\texttt 只用于 ASCII，避免再挂中文等宽字体触发 fontspec 查找
\usepackage[margin=2.6cm]{geometry}
\renewcommand{\abstractname}{摘要}
\renewcommand{\refname}{参考文献}
\renewcommand{\figurename}{图}
\renewcommand{\tablename}{表}

% ---- 数学与表格（须在 newtx 之前）----
\usepackage{amsmath,amssymb,amsthm}
\usepackage{booktabs,array,tabularx}
\usepackage[font=small,labelfont=bf,labelsep=quad]{caption}
\usepackage{fancyhdr}
\usepackage[numbers,sort&compress]{natbib}
\usepackage[hidelinks]{hyperref}

% ---- Times 风格（若已安装；务必放在 amssymb 之后，避免 \Bbbk 冲突）----
\IfFileExists{newtxtext.sty}{%
  \usepackage{newtxtext}%
  \IfFileExists{newtxmath.sty}{\usepackage{newtxmath}}{}%
}{}

\pagestyle{fancy}
\fancyhf{}
\fancyhead[L]{\small <短标题：不超过 60 字>}
\fancyhead[R]{\small \thepage}
\renewcommand{\headrulewidth}{0.4pt}
\setlength{\parskip}{2pt}
\setlength{\parindent}{2em}

% ---- 定理环境 ----
\theoremstyle{plain}
\newtheorem{theorem}{定理}[section]
\newtheorem{proposition}[theorem]{命题}
\newtheorem{lemma}[theorem]{引理}
\newtheorem{corollary}[theorem]{推论}
\theoremstyle{definition}
\newtheorem{definition}[theorem]{定义}
\renewcommand{\proofname}{证明}

\title{<标题>}
\author{Yu\\[2pt]{\small AI4Math 课程项目 · 研究性作业}}
\date{2026-10-23}

\begin{document}
\maketitle

\begin{abstract}
\noindent <摘要：问题是什么；本文做了什么；哪些结论由确定性工具认证（给出 claim\_id）；哪些仍是候选或未证。最后一句说明本文不宣称解决整个猜想。>

\medskip
\noindent\textbf{关键词}\quad <包含概率；Horvitz--Thompson 估计量；区组设计；bent 函数；极限分布>

\noindent\textbf{MSC 2020}\quad 62D05; 60F05; 05B05
\end{abstract}

\section{问题}
<猜想的精确陈述；量词顺序：$r$ 固定，$N\to\infty$，$n\to\infty$；已知结果。>

\section{构造}
<设计 $D_1$、$D_2$ 与总体。超平面
\begin{equation}\label{eq:hyperplane}
  H_{S,\varepsilon} = \{x \in \mathbb{F}_2^m : S\cdot x = \varepsilon\},
  \qquad S \neq 0,\ \varepsilon \in \{0,1\}.
\end{equation}
>

\section{命题与证明}
\begin{proposition}\label{prop:strength}
<对全部偶数 $m$，仿射超平面族是 $3$-$(N, N/2, N/4-1)$ 设计，且强度恰为 $3$。>
\end{proposition}
\begin{proof}
<计数证明；引用 \eqref{eq:hyperplane} 与账本里的陈述。>
\end{proof}

\section{精确验证}
\begin{table}[htbp]
  \centering
  \caption{有限实例的精确数值}
  \label{tab:instances}
  \begin{tabular}{@{}ccccccc@{}}
    \toprule
    $m$ & $N$ & $k$ & $b$ & $\lambda_3$ & D2 方差 & D2 到正态 KS \\
    \midrule
    4 & 16 & 8 & 30 & 3 & 16 & 0.195804 \\
    6 & 64 & 32 & 126 & 15 & 64 & 0.099328 \\
    8 & 256 & 128 & 510 & 63 & 256 & 0.05002 \\
    \bottomrule
  \end{tabular}
\end{table}

\section{局限与未证事项}
<逐条列出，与陈述账本里的 \texttt{unproved} 条目对应。>

\begin{thebibliography}{9}
\bibitem{chen2026} J. Chen. Low-order inclusion probabilities do not determine asymptotic normality. arXiv:2610.09168, 2026.
\bibitem{rothaus1976} O. S. Rothaus. On ``bent'' functions. \emph{Journal of Combinatorial Theory, Series A}, 20(3):300--305, 1976.
\bibitem{hajek1960} J. Hájek. Limiting distributions in simple random sampling from a finite population. \emph{Publ. Math. Inst. Hung. Acad. Sci.}, 5:361--374, 1960.
\bibitem{liding2017} X. Li, P. Ding. General forms of finite population central limit theorems with applications to causal inference. \emph{JASA}, 112(520):1759--1769, 2017.
\end{thebibliography}

\end{document}
